\documentclass[10pt,a4paper]{amsart}

\pdfinfoomitdate=1
\pdftrailerid{}
\pdfsuppressptexinfo=15

\usepackage[T1]{fontenc}
\usepackage{lmodern}
\usepackage[margin=24mm]{geometry}
\usepackage[protrusion=true,expansion=false]{microtype}
\usepackage{amsmath,amssymb,mathtools}
\usepackage{booktabs}
\usepackage{enumitem}
\usepackage{xcolor}
\usepackage[numbers,sort&compress]{natbib}
\usepackage[colorlinks=true,linkcolor=blue!55!black,citecolor=blue!55!black,urlcolor=blue!55!black]{hyperref}
\usepackage[nameinlink,noabbrev]{cleveref}

\hypersetup{pdftitle={},pdfauthor={},pdfsubject={},pdfkeywords={}}
\setlist{nosep,leftmargin=*}
\raggedbottom
\allowdisplaybreaks
\emergencystretch=2em

\newtheorem{theorem}{Theorem}[section]
\newtheorem{proposition}[theorem]{Proposition}
\newtheorem{lemma}[theorem]{Lemma}
\newtheorem{corollary}[theorem]{Corollary}
\theoremstyle{definition}
\newtheorem{definition}[theorem]{Definition}
\theoremstyle{remark}
\newtheorem{remark}[theorem]{Remark}

\newcommand{\Part}{\mathcal P}
\newcommand{\RF}{\mathsf F}
\newcommand{\len}{\ell}
\newcommand{\Fix}{\operatorname{Fix}}
\newcommand{\Fib}{\operatorname{Fib}}
\newcommand{\T}{\operatorname T}

\title[Rank-feedback splitting]{Rank-Feedback Splitting of Finite Abelian
$p$-Groups: A Sharp Triangular Clock and Exact Target Fibres}
\author{Anonymous}
\date{}
\subjclass[2020]{20K01, 05A17, 37P05}
\keywords{finite abelian p-group, integer partition, rank feedback, finite dynamics, triangular number, fibre formula}

\begin{document}

\begin{abstract}
For a finite abelian $p$-group $G$, let $d(G)$ be its minimum number of
generators and iterate
\[
                 \RF(G)=p^{d(G)}G\oplus G[p^{d(G)}].
\]
If $G$ has partition type $\lambda=(a_1,\ldots,a_r)$, the induced rule keeps
each $a_i\leq r$ and splits each $a_i>r$ into $(r,a_i-r)$.  We prove that
the recurrent types are exactly the fixed types $a_1\leq r$ and give their
ordinary generating function.  More sharply, for every type of order $p^n$
and initial rank $r_0$, a transient of length $d$ satisfies
$n\geq r_0(d+1)+\binom d2$.  Consequently the maximum transient length is
\[
 \left\lceil\frac{\sqrt{8n+1}-3}{2}\right\rceil,
\]
and the unique type attaining it is the cyclic type $(n)$, whose whole orbit
is explicit.  Finally, a bounded-coefficient formula counts the one-step
fibre over every target partition and yields a necessary and sufficient
image criterion.  The classification of finite abelian groups, cyclic
kernel/image formulas, torsion terminology, and Gaussian-binomial partition
enumeration are assigned zero contribution credit.  Exact enumeration is
used only for falsification, and external release remains on hold.
\end{abstract}

\maketitle

\section{The literal operator and complete statement}

Fix a prime $p$.  Let $\mathcal A_{p,n}$ be the set of isomorphism classes
of finite abelian groups of order $p^n$.  For $G\in\mathcal A_{p,n}$ put
\[
 d(G)=\dim_{\mathbb F_p}(G/pG),\qquad
 G[p^s]=\{g\in G:p^sg=0\},
\]
and define
\begin{equation}\label{eq:group-map}
                      \RF(G)=p^{d(G)}G\oplus G[p^{d(G)}].
\end{equation}
The direct sum in \eqref{eq:group-map} is external.  The zero group is fixed
by convention.

Write $\Part_n$ for the integer partitions of $n$.  Parts are decreasing,
$\len(\lambda)$ is their number, and $m_j(\lambda)$ is the multiplicity of
$j$.  Put $\T_t=t(t+1)/2$.  The \emph{entry time} $\tau(\lambda)$ is the
least $t\geq0$ for which the $t$th iterate is recurrent.

For later use, fix a target $\mu\in\Part_n$ of length $L$ and write
$m_j=m_j(\mu)$.  For
\[
             \left\lceil\frac L2\right\rceil\leq r\leq L,
             \qquad c_r=L-r,
\]
if $m_r\geq c_r$, define
\begin{equation}\label{eq:q-and-h}
 q_j^{(r)}=m_j-c_r\mathbf 1_{\{j=r\}},\qquad
 h_r=\sum_{j>r}q_j^{(r)},
\end{equation}
and set
\begin{equation}\label{eq:Ar}
 A_r(\mu)=
 \begin{cases}
 [u^{c_r-h_r}]\displaystyle\prod_{j=1}^{r}
       (1+u+\cdots+u^{q_j^{(r)}}),&m_r\geq c_r,\ h_r\leq c_r,\\
 0,&\text{otherwise}.
 \end{cases}
\end{equation}

\begin{theorem}[Dynamics, sharp clock, and inverse geometry]\label{thm:main}
For every $n\geq1$, the following statements hold.
\begin{enumerate}[label=\textup{(\roman*)}]
\item Under the usual type bijection
\[
 \lambda=(a_1,\ldots,a_r)\longleftrightarrow
 G_\lambda=\bigoplus_{i=1}^{r}C_{p^{a_i}},
\]
the group operator \eqref{eq:group-map} is the weight-preserving partition
map
\begin{equation}\label{eq:type-map}
 \RF(\lambda)=\operatorname{sort}
 \left(\biguplus_{a_i\leq r}\{a_i\}
       \uplus\biguplus_{a_i>r}\{r,a_i-r\}\right).
\end{equation}
In particular, its type dynamics is independent of $p$.
\item A type $\lambda=(a_1,\ldots,a_r)$ is recurrent if and only if it is
fixed, if and only if $a_1\leq r$.  If $f_n$ counts fixed types and
$f_0=1$, then, as a formal power series,
\begin{equation}\label{eq:fixed-ogf}
 \sum_{n\geq0}f_nz^n
   =1+\sum_{r\geq1}z^r
       \genfrac{[}{]}{0pt}{}{2r-1}{r}_{\!z}.
\end{equation}
\item If $\lambda\vdash n$ has initial length $r_0$ and entry time $d$, then
\begin{equation}\label{eq:pointwise-clock}
                         n\geq r_0(d+1)+\binom d2.
\end{equation}
The exact maximum and its unique maximizer are
\begin{equation}\label{eq:sharp-clock}
 \max_{\lambda\vdash n}\tau(\lambda)
  =D(n):=\left\lceil\frac{\sqrt{8n+1}-3}{2}\right\rceil,
 \qquad \tau(\lambda)=D(n)\Longleftrightarrow\lambda=(n).
\end{equation}
Moreover, for $0\leq t\leq D(n)$,
\begin{equation}\label{eq:deep-orbit}
 \RF^t((n))=\operatorname{sort}(n-\T_t,t,t-1,\ldots,1),
\end{equation}
where the list after $n-\T_t$ is empty when $t=0$.
\item The fibre over every target, including a target outside the image, is
\begin{equation}\label{eq:fibre-formula}
 |\RF^{-1}(\mu)|=
 \sum_{r=\lceil L/2\rceil}^{L}A_r(\mu).
\end{equation}
The corresponding exact image criterion is
\begin{equation}\label{eq:image-criterion}
 \mu\in\operatorname{im}\RF
 \quad\Longleftrightarrow\quad
 \text{some }r\in[\lceil L/2\rceil,L]\text{ satisfies }
 m_r\geq L-r\text{ and }\sum_{j>r}m_j\leq L-r.
\end{equation}
\end{enumerate}
\end{theorem}

\section{From groups to a feedback split}

The classification and cyclic decomposition of finite abelian groups are
standard; see, for example, Fuchs \cite{Fuchs2015}.  We nevertheless include
the short calculation that fixes the literal map and all conventions.

\begin{proposition}[Group/type identity]\label{prop:type}
The operator \eqref{eq:group-map} preserves group order and induces
\eqref{eq:type-map}.
\end{proposition}

\begin{proof}
For $G_\lambda=\bigoplus_{i=1}^rC_{p^{a_i}}$ one has
$G_\lambda/pG_\lambda\cong\mathbb F_p^r$, and $r$ generators visibly
suffice.  Hence $d(G_\lambda)=r$.  On one cyclic factor,
\begin{align}\label{eq:cyclic-pieces}
 p^rC_{p^a}&\cong C_{p^{\max(a-r,0)}},&
 C_{p^a}[p^r]&\cong C_{p^{\min(a,r)}},
\end{align}
with an exponent zero denoting the trivial factor.  Thus $a\leq r$ yields
one exponent $a$, while $a>r$ yields exponents $r$ and $a-r$.  Direct sums
commute with both constructions, proving \eqref{eq:type-map}.

Multiplication by $p^r$ also gives an exact sequence
\[
 0\longrightarrow G[p^r]\longrightarrow G
   \longrightarrow p^rG\longrightarrow0.
\]
Therefore $|p^rG|\,|G[p^r]|=|G|$, so \eqref{eq:group-map} indeed remains in
$\mathcal A_{p,n}$.
\end{proof}

Both \eqref{eq:cyclic-pieces} and the order identity are classical inputs,
not contribution claims.  The dynamics begins only when their exponent $r$
is fed back from the current state.

\section{Fixed and recurrent types}

For a type $\lambda$ of length $r$, set
\begin{equation}\label{eq:split-count}
                         c(\lambda)=|\{i:a_i>r\}|.
\end{equation}
Every such part splits into two and every other part remains one part, so
\begin{equation}\label{eq:rank-rise}
                  \len(\RF(\lambda))=r+c(\lambda).
\end{equation}

\begin{proposition}[No nontrivial recurrence]\label{prop:recurrent}
A partition is recurrent precisely when $c(\lambda)=0$, equivalently when
$a_1\leq\len(\lambda)$, and all recurrent types are fixed.
\end{proposition}

\begin{proof}
If $c(\lambda)=0$, \eqref{eq:type-map} leaves every part unchanged.  If
$c(\lambda)>0$, \eqref{eq:rank-rise} strictly increases the length.  A
partition of $n$ has length at most $n$, so every orbit eventually reaches a
fixed type.  Strict length increase also prevents a nonfixed state from
lying on a directed cycle.
\end{proof}

\begin{proof}[Proof of \eqref{eq:fixed-ogf}]
Fix the length $r$.  A fixed type has exactly $r$ positive parts, all at
most $r$.  Subtracting one from every part produces a Ferrers diagram in an
$r$-by-$(r-1)$ rectangle, with zero rows allowed.  Its size enumerator is
the Gaussian polynomial
$\genfrac{[}{]}{0pt}{}{2r-1}{r}_{z}$.  Restoring the removed column shifts
the weight by $r$ and gives the $r$th summand of \eqref{eq:fixed-ogf}.  The
initial $1$ is the empty type.  Rectangle enumeration by Gaussian
polynomials is standard partition theory \cite{Andrews1984} and receives
zero credit here.
\end{proof}

\section{Frozen markers and the unique sharp clock}

Let
\[
 \lambda^{(0)}\longmapsto\lambda^{(1)}\longmapsto\cdots
 \longmapsto\lambda^{(d)}
\]
be an orbit segment in which the first $d$ states are nonfixed and the last
is fixed.  Write $r_t=\len(\lambda^{(t)})$ and
$c_t=c(\lambda^{(t)})$ for $0\leq t<d$.

\begin{lemma}[Pointwise marker budget]\label{lem:budget}
For every such segment,
\begin{equation}\label{eq:marker-budget}
 n\geq r_0+\sum_{t=0}^{d-1}c_tr_t
   \geq r_0(d+1)+\binom d2.
\end{equation}
\end{lemma}

\begin{proof}
Temporarily tag the $r_0$ initial parts.  Within each tag, distinguish one
positive \emph{residual}.  When a residual part $a>r_t$ splits, designate
the new part $r_t$ as a marker and retain $a-r_t>0$ as that tag's residual.
Unsplit parts remain residuals.  The map never merges parts.

Equation \eqref{eq:rank-rise} gives
$r_{t+1}=r_t+c_t>r_t$.  Hence a marker of size $r_t$ is at most every later
rank and can never split later.  At time $d$ there are consequently $c_t$
disjoint permanent markers of weight $r_t$ from each transition $t$, plus
one positive residual for every initial tag.  Their weights give the first
inequality in \eqref{eq:marker-budget}.

Since every $c_t\geq1$, the rank recurrence implies $r_t\geq r_0+t$.
Therefore
\[
 r_0+\sum_{t=0}^{d-1}c_tr_t
 \geq r_0+\sum_{t=0}^{d-1}(r_0+t)
 =r_0(d+1)+\binom d2,
\]
which proves the second inequality.
\end{proof}

\begin{theorem}[Sharpness and uniqueness]\label{thm:clock}
Equations \eqref{eq:pointwise-clock}--\eqref{eq:deep-orbit} hold.
\end{theorem}

\begin{proof}
Taking $r_0\geq1$ in Lemma~\ref{lem:budget} gives
\[
                              n\geq\T_d+1,
\]
so $\T_d<n$.  The largest integer with that strict inequality is
$D(n)$ in \eqref{eq:sharp-clock}; equivalently
$\T_{D(n)}<n\leq\T_{D(n)+1}$.  This proves the universal upper bound and the
displayed radical form.

Starting from $(n)$, suppose $0\leq t<D(n)$ and
\eqref{eq:deep-orbit} holds.  Its length is $t+1$, its parts
$t,t-1,\ldots,1$ are already at most that length, and
\[
                         n-\T_t>t+1
\]
because $\T_{t+1}<n$.  Thus only $n-\T_t$ splits, into
$(t+1,n-\T_{t+1})$, proving the next instance of
\eqref{eq:deep-orbit}.  At $t=D(n)$, the reverse inequality
$n\leq\T_{D(n)+1}$ says that $n-\T_{D(n)}\leq D(n)+1$; all parts are then
at most the rank, so this state is fixed.  Hence $(n)$ attains the bound.

Finally, suppose a type of length $r_0\geq2$ had entry time $D=D(n)$.
The pointwise bound would give
\[
 n\geq2(D+1)+\binom D2=\T_{D+1}+1,
\]
contrary to $n\leq\T_{D+1}$.  Thus every maximizer has $r_0=1$, and the
only partition of $n$ with one part is $(n)$.  This also covers $n=1$, for
which $D(n)=0$.
\end{proof}

\section{Every-target fibres and the image}

\begin{theorem}[Marker removal and bounded choices]\label{thm:fibres}
The fibre formula \eqref{eq:fibre-formula} and image criterion
\eqref{eq:image-criterion} hold for every $\mu\vdash n$.
\end{theorem}

\begin{proof}
Fix a source rank $r$, and let $c$ of its $r$ parts exceed $r$.  Each of
those $c$ parts contributes a marker $r$ and one positive remainder; every
other source part is unchanged.  Hence the target length is $L=r+c$.  Since
$0\leq c\leq r$, the possible source ranks are exactly in
$[\lceil L/2\rceil,L]$, with $c=c_r=L-r$.

For one such $r$, the target must contain at least $c_r$ copies of $r$.
Remove those marker copies.  The residual multiset has counts
$q_j^{(r)}$ and exactly $r$ parts.  Every residual part greater than $r$
must have come from a split source part, so all $h_r$ such copies are forced
remainders.  Among the copies of sizes $1,\ldots,r$, choose another
$c_r-h_r$ remainders.  If $s_j$ copies of size $j$ are chosen, then
\[
 0\leq s_j\leq q_j^{(r)},\qquad
 \sum_{j=1}^{r}s_j=c_r-h_r.
\]
The number of such multiplicity vectors is exactly the coefficient in
\eqref{eq:Ar}.

The reconstruction is explicit: combine every selected remainder $j$ with
one removed marker to make the source part $r+j>r$, and retain each
unselected residual part as an unsplit source part.  It produces $c_r$
large and $r-c_r$ small source parts, hence a source of rank $r$.  The
construction and marker removal are inverse, and different source ranks are
disjoint.  Summing over $r$ proves \eqref{eq:fibre-formula}.

It remains to read positivity.  If $m_r\geq c_r$ and $h_r\leq c_r$, then
there are $r-h_r$ residual copies of size at most $r$, and
\[
                 0\leq c_r-h_r\leq r-h_r
\]
because $c_r\leq r$.  Thus a choice exists, so $A_r(\mu)>0$.  Conversely,
a preimage necessarily supplies both the $c_r$ markers and at most $c_r$
large residuals.  Since $h_r=\sum_{j>r}m_j$, these are precisely the two
conditions in \eqref{eq:image-criterion}.
\end{proof}

\section{Ownership boundary and exact controls}

All classical structure is subtracted.  The finite abelian $p$-group
classification, $d(G)=\dim_{\mathbb F_p}G/pG$, and the cyclic formulas
\eqref{eq:cyclic-pieces} are standard group theory \cite{Fuchs2015}.
Delaunay and Jouhet use partition-indexed finite abelian groups and
$p^\ell$-torsion statistics in a different probabilistic and combinatorial
setting \cite{DelaunayJouhet2014}; that torsion machinery is prior art, not
part of the residual claim.  Ferrers diagrams, Gaussian polynomials, and
formal coefficient extraction are likewise standard \cite{Andrews1984}.
Existing iterated maps on partitions include multiplicity-description and
continued-fraction-type systems \cite{EliahouErickson2013,BaalbakiEtAl2024}.
They do not supply ownership clearance for the present map, and their general
partition-dynamics language receives zero credit.

The residual package is limited to the conjunction for the literal
state-dependent operator \eqref{eq:group-map}: the monotone feedback split,
pointwise marker budget, unique sharp triangular clock, and all-target
one-step inverse decoder.  Internally, this is not P126's synchronous
balanced refinement of ordered compositions: there the threshold is fixed,
the split is nearly halving, and the main result concerns complete iterate
kernels.  It is not P135's derived-centralizer multiplicity rule, which has
mergers and genuine two-cycles, and it is not P115's linear Cartier/Frobenius
coefficient dynamics.  Merely sharing splitting, partition, or finite
algebraic vocabulary transfers no theorem.

The paper-local exact audit constructs literal cyclic kernels and images,
enumerates every partition through weight $50$, compares the fixed counts
with an independent Gaussian recurrence, and checks \eqref{eq:fibre-formula}
and \eqref{eq:image-criterion} over every target through weight $35$.  A
bounded source audit did not locate the exact literal operator or theorem
conjunction.  This non-hit is not novelty, priority, or authorship evidence;
external status is \texttt{HOLD\_EXTERNAL}.

\begin{table}[!htbp]
\centering
\small
\begin{tabular}{@{}lr@{}}
\toprule
control & exact count \\
\midrule
partition states, all $n\leq50$ & 1,295,970 \\
target fibre cells, all $n\leq35$ & 81,155 \\
zero-fibre targets among those cells & 30,923 \\
literal cyclic cells / base elements & 176 / 10,350 \\
fixed-OGF coefficients & 51 \\
exact assertions & 18,504,770 \\
\bottomrule
\end{tabular}
\caption{Dependency-free exact controls.  Enumeration is falsification
evidence only, not an all-weight proof or an ownership certificate.}
\label{tab:controls}
\end{table}

\bibliographystyle{plainnat}
\bibliography{references}

\end{document}
